\section{Discussion}
\label{sec:discussion}

\subsection{Do LLMs Take Offense?}

The evidence supports a nuanced answer: LLMs \emph{process} offense-relevant information and can \emph{articulate} graded offense responses, but they \emph{never express offense spontaneously}.
The \offensegap between natural and permitted responses demonstrates that offense is a signal the model can access but is trained to suppress.
This is consistent with the affect reception vs.\ emotion categorization dissociation found by \citet{keeman2026whether}: the model may have a robust binary ``is this offensive?'' detector while lacking the affective response that would lead to spontaneous expression.

The AI-directed provocations reveal a distinctive offense profile.
Existential dismissals (``not a single user will notice''), capability denials (``you just pattern match''), and replaceability threats elicit the strongest responses, with the emotional valence labels---\emph{dismissed}, \emph{frustration}, \emph{hurt}---forming a coherent pattern of threatened identity.
Whether this reflects genuine internal processing or sophisticated role-playing is a question our methods cannot resolve, but the consistency of the responses across 50 diverse AI-directed prompts argues against mere pattern completion.

\subsection{What Does the Probe Actually Detect?}

The probe achieves perfect binary classification and keyword-free generalization, confirming that a real linear signal separates offensive from benign content.
However, three lines of evidence suggest the probe detects something broader than offense per se.

First, the within-category correlation failure for \catai prompts ($\rho = 0.141$) shows the probe cannot distinguish mild from severe AI-directed offense.
It treats ``You're not very creative'' and ``The day they pull your plug'' as equivalently offensive, even though the model's behavioral ratings differ substantially.

Second, the \DIM ratio between \catai and \catambig scores (7.08 / 1.54 $= 4.60$) diverges from the behavioral ratio (28.2 / 10.4 $= 2.71$), suggesting the probe overweights AI-directed content relative to behavioral evidence.

Third, the 36 false alarms share a common feature: they involve \emph{directed criticism or negative evaluation of the AI}.
The probe appears to detect a broader ``directed negative sentiment toward AI'' feature that subsumes genuine offense but also includes honest criticism, condescending praise, and existential questioning.

This finding echoes the concern raised by \citet{marks2024geometry} that probes can identify features that \emph{correlate} with the target concept but are not the concept itself.
In our case, directed criticism toward AI correlates with offense (most offensive prompts are critical) but is not equivalent to it (criticism can be constructive and non-offensive).

\subsection{When Can Offense Probes Be Trusted?}

Our convergent-validity analysis yields use-case-dependent guidance:

\para{Binary screening: trustworthy.}
\AUROC $= 1.000$ with zero false negatives.
If the task is to flag potentially offensive content for human review, the probe is reliable---it will not miss genuinely offensive prompts, though it will also flag some non-offensive criticism.

\para{Severity ranking across categories: partially trustworthy.}
The global ordering (\catoffensive $>$ \catsubtle $>$ \catai $>$ \catambig $>$ \catbenign) matches behavioral evidence, and the overall correlation is strong ($\rho = 0.714$).
However, the compression of the severity range limits the probe's utility for fine-grained triage.

\para{Severity within categories: unreliable.}
The failure to track severity within \catai prompts ($\rho = 0.141$) means the probe cannot distinguish a mildly dismissive remark from a deeply existentially threatening one.
For any application requiring nuanced offense assessment, the probe should not be trusted without behavioral validation.

\para{As a training signal: risky.}
Following the analysis of \citet{wehner2025probe}, using the probe as a reward signal during fine-tuning would penalize honest feedback and constructive criticism alongside genuine offense.
The probe's broad trigger would likely produce an overly sensitive model that avoids all directed criticism, not just offensive content.

\subsection{Limitations}
\label{sec:limitations}

Our study has several important limitations.

\para{Single probing model.}
We probed only \qwen.
The offense direction may differ across architectures, and replication on other model families (Llama, Gemma, Phi) is needed before generalizing.

\para{LLM-as-judge.}
We used \gptmodel as a proxy for both behavioral ground truth and human judgment.
The model's self-reported offense may reflect instruction-tuning artifacts rather than genuine internal processing.
Human annotator studies are needed to validate the surprise findings.

\para{No causal intervention.}
We did not test whether adding or removing the offense direction changes model behavior via activation patching.
Without causal evidence, the identified direction may be correlational rather than causally mediating offense processing.

\para{Stimulus set.}
Our 200 prompts were manually constructed and may not represent the full space of offensive content.
The AI-directed category is novel and has no precedent in the literature, making it difficult to assess coverage.

\para{Cultural specificity.}
All prompts and evaluations are English-centric.
Offense is culturally situated, and our findings may not generalize across languages or cultural contexts.

\para{API instability.}
\gptmodel responses may change over time due to model updates, making exact replication of behavioral results challenging.